\documentclass[11pt]{article}

\usepackage[margin=1in]{geometry}
\usepackage[T1]{fontenc}
\usepackage[utf8]{inputenc}
\usepackage{microtype}
\usepackage{hyperref}
\usepackage{enumitem}
\usepackage{booktabs}
\usepackage{amsmath}

\hypersetup{
  colorlinks=true,
  linkcolor=black,
  citecolor=black,
  urlcolor=black
}

\title{Observer-Declared Timelines for Real-World Agent Work}
\author{Keren Dong\\kungfu-systems}
\date{Draft: July 2026}

\begin{document}

\maketitle

\begin{abstract}
Software systems that manage real-world agent work increasingly combine facts
from local machines, remote sessions, repositories, user interfaces, command
line tools, model providers, and human decisions. These systems often present a
single timeline even though the underlying facts come from multiple clocks,
locations, and participants. This paper argues that such timelines should not
be presented as views from nowhere. A useful timeline should declare its
observer, accepted facts, projection policy, causal constraints, and degraded
evidence state. We call this design rule \emph{observer-declared timelines}.
The rule is motivated by KFD-4, a kungfu-systems procedure derived from the
foundation that facts must not drift, trust must start from facts, and
cooperation must start from trusted value. We define the problem, give a
minimal model, contrast the approach with global clocks, central sequencers,
event sourcing, distributed tracing, and CRDTs, and outline how Kungfu can use
the rule as part of a runtime fact ledger for real-world agent work.
\end{abstract}

\section{Introduction}

Consider a long-running agent task that crosses two sessions. The first
session changes a repository, starts remote work, records one failed check, and
ends before the task is complete. A later participant opens the local product
and asks what happened and what action is safe now. The local interface, remote
runtime, repository, build system, and retained session evidence may each
contain valid facts. They need not contain the same facts, reach the same cut,
or order concurrent facts in the same way.

The KFD Foundation Triad describes how such work can remain one shared world
across participants and time \cite{kfdFoundation2026}:

\begin{equation}
\text{Facts} \rightarrow \text{Trust} \rightarrow \text{Cooperation}
\rightarrow \text{New Facts}.
\end{equation}

KFD-1 preserves stable reference, KFD-2 binds reliance to evidence, and KFD-3
makes relevant value and constraints visible to a participant before action.
Cooperation then creates consequences that return through review into a
successor fact world. That loop solves a continuity problem, but it creates a
new pressure: no participant consumes the entire world from every position.

The pressure appears separately at each layer of the triad. If a selected or
transformed fact cut is presented without perspective, selection can masquerade
as reality and violate the stable reference required by KFD-1. If a claim does
not identify which cut and projection support it, a later participant cannot
reconstruct the warrant required by KFD-2. If a view hides whose consequence
position, available actions, and known gaps it represents, a participant cannot
inspect the value and constraints needed by KFD-3; cooperation degrades into
guessing, obedience, or repeated reconstruction.

This yields the first derived pressure from the Foundation Triad:

\begin{equation}
\text{non-drifting facts} + \text{multiple participants}
\Longrightarrow \text{declared perspective}.
\end{equation}

KFD-4 names the corresponding rule: timelines must declare their observer
\cite{kfd2026observer}. An observer-declared timeline does not deny facts or
causality. It states which observer position is represented, which facts were
accepted, how the view was projected, which causal relations dominate that
projection, what consequences or actions the view is meant to support, and
where evidence remains incomplete. The result is not a view from nowhere. It
is a reproducible claim about how one view was formed from a shared fact world.

The core rule leads to a second question only after the perspective is defined:
what happens when the observer implementation itself restarts, upgrades, or is
replaced? A stable product label does not prove that a new process accepted the
same facts or applied the same projection semantics. We therefore distinguish a
logical observer from its incarnations and treat continuity between them as an
evidence-bearing transition. This is a derived refinement of the core rule, not
a prerequisite for first understanding why perspective must be declared.

This paper makes five contributions:

\begin{itemize}[leftmargin=*]
  \item a derivation of observer declaration from the fact, trust, and
        cooperation responsibilities of the KFD Foundation Triad;
  \item a minimal model over accepted fact cuts, causal constraints, projection
        policy, consequence position, degradation, and verification;
  \item an incarnation and continuity-witness refinement for observers that
        fail, restart, or change implementation;
  \item a mapping from that generic model to Kungfu's Episode fact authority,
        exact-cut readiness, fenced recovery, and product surfaces; and
  \item bounded implementation evidence together with an explicit account of
        the observer-level evaluation that remains incomplete.
\end{itemize}

Section~\ref{sec:background} develops the problem through the running example
and distributed-systems context. Section~\ref{sec:model} defines the generic
model without requiring a Kungfu-specific causal object. Section~\ref{sec:design}
then introduces Episodes and runtime mechanisms as one concrete realization.
Section~\ref{sec:evaluation} separates evidence for that substrate from the
still-open end-to-end observer claim, and Section~\ref{sec:related} compares
related work.

\section{Background and Motivation}
\label{sec:background}

\subsection{Real-world agent work}

We use \emph{real-world agent work} to mean work in which agents help users
manage concrete tasks rather than only produce isolated model outputs. Such
work may involve local state, external repositories, remote terminals, GUI and
TUI controls, model calls, package releases, documents, and human decisions.
The product problem is not only how to run an agent. It is how to keep the
state of the work legible over time.

Long-lived work introduces a persistence problem. Users and agents need to ask:
what happened, what source observed it, what evidence supports it, what is
missing, and what action is safe now? A timeline view is a natural interface
for these questions, but only if the view makes its own assumptions inspectable.

\subsection{The global-clock temptation}

Distributed systems theory has long shown that time and ordering are not
trivial in distributed computation. Lamport's happened-before relation
separates causal order from physical clock order \cite{lamport1978time};
vector-clock style systems make concurrency explicit \cite{mattern1989virtual}.
Yet many product surfaces still sort mixed-source facts primarily by wall-clock
timestamps and present the result as a single history. That is often sufficient
for casual logs, but it is not sufficient for recovery, audit, or cooperation
among capable participants.

The problem is sharper in agent systems because the user may not have the time
or expertise to manually diagnose missing evidence. A local product must help
both the human and the agent understand what the view means.

\subsection{KFD foundation}

The KFD foundation used by Kungfu can be summarized as:

\begin{quote}
KFD-1: facts must not drift.\\
KFD-2: trust must start from facts.\\
KFD-3: cooperation must start from trusted value.
\end{quote}

KFD-4 applies this foundation to timelines:

\begin{quote}
KFD-4: timelines must declare their observer.
\end{quote}

Under this framing, a timeline view is trustworthy only when it is grounded in
non-drifting facts, exposes enough evidence for trust assessment, and helps
participants cooperate through transparent value rather than hidden pressure or
opaque authority \cite{kfd2026observer}.

\subsection{Why a causal object is needed}

A projection rule alone does not define the unit that can be accepted, moved,
verified, or rejected. Raw journal pages are physical append blocks; sessions
and runs are product vocabulary; individual frames are too small to provide a
portable closure boundary. Kungfu therefore introduces an \emph{Episode}: a
bounded causal segment containing the manifest, frame references, content
commitments, provenance, schemas, dependencies, and verification evidence
needed to inspect one unit of work \cite{kungfu2026episode}.

Episodes do not make the timeline absolute. They make its inputs explicit. A
view can name which Episodes it accepted, preserve their internal and declared
cross-Episode causality, and then apply observer policy only where the causal
facts leave an ordering choice.

\section{Observer-Declared Timeline Model}
\label{sec:model}

An observer-declared timeline is a deterministic projection over accepted
facts from a declared perspective. The model is intentionally small; it is a
surface contract, not a full database schema.

\subsection{Facts and sources}

A \emph{fact} is an event, record, manifest, receipt, payload hash, causal
edge, or verification result that the system treats as evidence. A fact has a
source, provenance, and a local identity. A fact source may be a local runtime,
remote runtime, repository, release artifact, adapter, imported bundle, or
external system.

The view does not need to accept every fact the product has ever seen. It must
declare which source ranges, heads, cursors, manifests, watermarks, or
freshness boundaries were accepted for the current projection.

\subsection{Observer}

An \emph{observer} is the perspective whose view is being shown. It may be a
human, agent, runtime location, service, or product view. The observer is not
necessarily the author of the facts. It is the subject from whose accepted
fact set and projection policy the timeline is derived.

\subsection{Projection policy}

A projection policy orders accepted facts for display, replay, export, or
agent consumption. Causal constraints dominate projection policy: if fact
$a$ is known to causally precede fact $b$, a valid projection must not place
$b$ before $a$. Concurrent or causally unrelated facts may be ordered by an
observer policy, source priority, local sequence, or deterministic
tie-breaker, but the policy must be inspectable.

\subsection{Degraded evidence}

A timeline may be useful even when evidence is incomplete. Instead of hiding
that incompleteness, the view should declare degraded states such as missing
causality, incomplete accepted ranges, stale remote facts, redacted payloads,
schema mismatch, unverifiable source, or unresolved conflict. A degraded view
is preferable to an overconfident total order.

\subsection{Minimal surface}

The minimal observer-declared surface contains:

\begin{itemize}[leftmargin=*]
  \item a stable view identifier;
  \item the view subject, such as timeline, history, replay, sync, ordering,
        or mixed-source work state;
  \item observer identity and kind;
  \item accepted fact sources and ranges;
  \item projection policy version, causal dominance, and tie-breaker;
  \item causal constraints;
  \item degraded evidence states; and
  \item verification status.
\end{itemize}

This is close to the KFD-4 observer-perspective schema published by the KFD
package. The schema is a vocabulary and verification boundary, not a claim that
an adopter's runtime implementation is correct merely because it imports the
schema.

\section{Kungfu Prototype and Design}
\label{sec:design}

Kungfu is infrastructure for managing real-world work with agents. In this
context, an observer-declared timeline is one surface of a broader runtime fact
ledger. The current implementation establishes the Episode substrate beneath
that surface; it does not yet implement the complete multi-machine observer
projection.

\subsection{Runtime fact ledger}

The runtime fact ledger preserves source-local facts, provenance, causal links,
payload evidence, schema bindings, and verification results. Future sync views
also need accepted source ranges and freshness boundaries. The visible timeline
is a projection over the ledger, not the ledger itself.

This distinction matters. If storage flattens remote facts into anonymous local
records, the product loses the ability to explain what was accepted and why. If
the GUI presents a global-looking order without observer metadata, the user or
agent cannot reproduce the view during recovery.

\subsection{Implemented Episode path}

The Episode prototype separates five layers \cite{kungfu2026episode}:

\begin{enumerate}[leftmargin=*]
  \item An append-only yijinjing manifest journal stores fixed-layout Episode
        lifecycle, frame, reference, and dependency records as authority.
  \item One typed C++ fold derives the canonical current Episode view. Python,
        Node, CLI, and GUI edges do not define independent Episode semantics.
  \item Content-addressed immutable storage resolves payload and schema
        references, with hashes and lengths verified before use.
  \item Structural fsck checks manifest integrity, causal closure, dependencies,
        content evidence, lifecycle state, and projection drift. Recovery and
        repair append evidence instead of silently rewriting verified history.
  \item A SQLite Episode projection accelerates query but remains rebuildable
        from the journal. Absence or drift degrades the projection without
        promoting the database to authority.
\end{enumerate}

This architecture provides the accepted-object and evidence substrate required
by an observer-declared timeline. It also makes partial failure visible at the
Episode boundary: lifecycle, health, and operation capability are separate
dimensions.

\subsection{KFD-4 implementation map}

Table~\ref{tab:kfd4-map} distinguishes implemented evidence from the remaining
observer-view work.

\begin{table}[ht]
\centering
\small
\begin{tabular}{>{\raggedright\arraybackslash}p{0.22\linewidth}>{\raggedright\arraybackslash}p{0.34\linewidth}>{\raggedright\arraybackslash}p{0.34\linewidth}}
\toprule
KFD-4 field & Current substrate & Remaining view work \\
\midrule
Observer & Source and writer-location provenance & First-class observer and stable view identifier \\
Accepted facts & Episode manifests, refs, dependencies, export/import evidence & Multi-source accepted ranges, watermarks, and freshness \\
Projection policy & Deterministic typed fold and rebuildable Episode query projection & Versioned cross-source ordering policy and tie-breaker \\
Causal constraints & In-Episode closure and declared Episode dependencies & Projection-level enforcement over a selected multi-source set \\
Degraded evidence & Structured fsck issues, missing evidence, stale projection, recovery paths & View-level degradation for incomplete ranges and policy knowledge \\
Verification & Semantic oracle, fsck, Trust Report, CI gate & Exported observer declaration and projection-level fsck \\
\bottomrule
\end{tabular}
\caption{Current Kungfu evidence mapped to the KFD-4 observer-perspective contract.}
\label{tab:kfd4-map}
\end{table}

\subsection{Dual-first interfaces}

The same timeline contract should be usable by humans and agents. A GUI can
show the default observer view while exposing metadata on demand. A CLI or API
can return the same accepted facts, projection policy, and degraded evidence
state directly. Agent-facing surfaces are not secondary integrations; they are
first-class control surfaces for participants that may inspect the system more
frequently than humans.

\subsection{Remote work}

Real-world agent work often happens on remote machines. A server-side agent may
write facts into a remote Kungfu home while the local GUI remains the user's
primary control plane. A useful product should support importing or syncing
remote facts back to the local perspective without pretending that the local
observer saw every event live. The imported timeline should declare source
location, accepted range, freshness, and any degraded evidence.

\subsection{Verification gates}

KFD-4 can be used as a release or product gate for perspective-bearing
surfaces. Passing the gate should mean that the surface exposes the observer,
accepted facts, projection policy, causal constraints, and degraded evidence
metadata. It should not mean that every adapter, remote source, or runtime
payload is complete. Stronger claims require their own KFD-2-style trust
assessment.

The Episode qualification report is an example of that evidence discipline. It
records source revision, profile, workload, semantic dimensions, violations,
performance observations, and known gaps. It is evidence for a bounded Episode
profile, not an automatic KFD-4 certificate for every timeline surface.

\section{Evaluation Plan}
\label{sec:evaluation}

This draft does not yet report empirical results. The planned evaluation should
test whether observer-declared timelines improve recovery, auditability, and
agent cooperation in realistic workflows.

\subsection{Case studies}

\begin{itemize}[leftmargin=*]
  \item \textbf{Remote agent work.} An agent performs a task on a server while
        the user later reviews the work from a local GUI.
  \item \textbf{Multi-machine sync.} Facts from local and remote runtimes are
        imported with different freshness boundaries and causal completeness.
  \item \textbf{Replay and debugging.} A user or agent reconstructs why a task
        reached a given state after partial failure.
  \item \textbf{Release evidence.} Build and release facts are shown as a
        timeline with source provenance and verification status.
\end{itemize}

\subsection{Measures}

Potential measures include:

\begin{itemize}[leftmargin=*]
  \item recovery time after interruption;
  \item number of hidden-ordering assumptions discovered during audit;
  \item ability for an agent to reproduce a displayed timeline from exported
        metadata;
  \item correctness of causal-order preservation under clock skew;
  \item rate of degraded-evidence states surfaced rather than silently hidden;
        and
  \item user and agent ability to decide the next safe action from the view.
\end{itemize}

\subsection{Artifacts}

Evaluation artifacts should include fixture ledgers, projected timeline
exports, KFD-4 observer-perspective declarations, and KFD-2 trust assessments.
The paper should not rely on private operational logs for evidence unless they
are sanitized, permissioned, and explicitly marked.

\section{Related Work}
\label{sec:related}

\subsection{Logical and physical time}

Lamport clocks and happened-before relations establish that causal order is not
the same as physical timestamp order \cite{lamport1978time}. Vector-clock
systems make concurrent events explicit \cite{mattern1989virtual}. The
observer-declared timeline rule does not replace these mechanisms. It brings a
similar discipline to product surfaces: when the product shows a timeline, it
should expose enough perspective metadata for the ordering to be audited.

Distributed snapshots construct a consistent global state without requiring a
physically simultaneous observation \cite{chandy1985snapshots}. An
observer-declared view has a related concern, but its contract extends beyond
one snapshot algorithm: it names accepted heterogeneous facts, projection
semantics, evidence degradation, and the incarnation that claims the view.

\subsection{Failure detection and recovery epochs}

Failure detectors reason about what distributed processes can suspect and what
reliable services can be built from those suspicions \cite{chandra1996failure}.
The incarnation model addresses a product-layer consequence: replacement must
not turn liveness evidence into authority continuity. Generation and epoch
fences are implementation techniques for rejecting stale actors; the proposed
surface additionally asks a timeline to expose the cut and semantic relation
that make a successor view reproducible.

\subsection{Event sourcing}

Event sourcing treats state as a derivation from an append-only sequence of
events \cite{fowler2005eventsourcing}. Observer-declared timelines share the
idea that records matter, but differ in emphasis. The visible timeline is not
assumed to be the authoritative event log. It is a perspective-bearing
projection over accepted facts that may come from several sources. Episodes add
a bounded causal and verification object above physical append blocks, while
keeping projections rebuildable from authority.

\subsection{Distributed tracing}

Distributed tracing systems such as OpenTelemetry capture spans and causal
relationships across services \cite{opentelemetry}. This is valuable evidence
for runtime behavior. The proposed rule is broader at the product surface: a
timeline may combine traces, repository facts, user decisions, release
manifests, local files, and remote agent actions. An Episode may contain trace
evidence, but its boundary is defined by causal closure and portable trust rather
than by one tracing backend or span vocabulary.

\subsection{CRDTs and replicated state}

CRDTs provide convergence properties for replicated data types under
concurrent updates \cite{shapiro2011crdt}. Observer-declared timelines are not
a conflict-free data structure. They are a view contract: even when underlying
state uses CRDTs, consensus, or a central sequencer, a mixed-source product
timeline should still declare what perspective the user or agent is seeing.

\subsection{Human-computer and human-agent cooperation}

Engelbart framed computing as augmenting human intellect rather than merely
automating tasks \cite{engelbart1962augmenting}. Real-world agent work extends
that problem. The product must support cooperation between human and agent
participants that observe, remember, and act through different surfaces.

\section{Discussion and Limitations}
\label{sec:discussion}

\subsection{Not relativism}

Declaring an observer does not mean that any story is acceptable. Facts,
causality, hashes, manifests, schemas, and verification results remain
load-bearing. The rule rejects an undeclared view from nowhere, not the
existence of facts.

\subsection{Not universal consensus}

Some systems require consensus, total ordering, or a central sequencer. The
observer-declared timeline rule does not forbid those designs. It says that
when a product presents a timeline to a human or agent, the view should still
declare the accepted facts and projection policy behind it.

\subsection{Adopter responsibility}

A schema can standardize vocabulary, but it cannot prove an adopter's runtime
correctness by itself. An adopter that claims a concrete timeline is complete,
fresh, or causally valid must provide product-specific evidence. This is why
KFD-4 should be assessed through KFD-2-style trust claims when used as a
release or product gate.

\subsection{Research status}

This is an initial design paper. The next version should replace design
intentions with concrete Kungfu implementation evidence: exported observer
perspectives, replay fixtures, failure cases, and measured recovery outcomes.

\section{Conclusion}

Real-world agent work makes timeline views more important and more dangerous.
They are important because humans and agents need durable surfaces for recovery,
audit, and cooperation. They are dangerous when they hide their perspective and
present mixed-source facts as an absolute global order.

Observer-declared timelines offer a practical alternative. A product can keep
facts and causality load-bearing while admitting that every useful view is
projected from a declared perspective. This paper adds a further constraint:
the observer is not timeless. A stable logical observer may have many
incarnations, and continuity between them must be established through accepted
cuts, policy/runtime identity, recovery evidence, and explicit degradation.

Kungfu now supplies preliminary implementation evidence for both sides of this
rule. Episodes provide bounded causal authority and safe-capability evidence.
Exact-cut runtime readiness, generation and coordinator-epoch fencing, Peer
rebootstrap, explicit recovery states, and immutable runtime selection show how
an observer incarnation can fail or change without silently acquiring old
authority. Historical scale evidence and newer source-bound single-host
qualification remain separate, preserving the provenance of each claim.

The results do not yet validate the complete observer-relative multi-machine
view. That next step requires first-class observer policy, accepted source
ranges, reproducible cross-source ordering, exported incarnation and continuity
declarations, projection-level causal fsck, and experiments in which humans and
agents recover from the same declared perspective. Keeping that boundary
explicit is part of the proposal itself: trust in a timeline, including trust in
its continuity, must start from the facts actually available.


\bibliographystyle{plain}
\bibliography{references}

\end{document}
